\emph{Setting.} Work in units of $\kappa$. The array of ratio $r$ and twist $\delta$ has modes $k \ge 0$ of weight $(k + \delta)/r$ above the sea and, for $k < 0$, holes of weight $(k' + 1 - \delta)/r$, $k' = -1 - k \ge 0$. Write $W$ for the set of numbers $k + \delta$ and $k + 1 - \delta$, $k \ge 0$. By Theorem \claimref{synthesis.null-ordering-of-counted-units-is-zeta-symmetric}(ii), applied to the spacing $\kappa/r$ in place of $\kappa$ (the zeta-regularised sum is homogeneous of degree one in the spacing), the sea energy is $B_{2}(\delta)/(2r)$. In the Gibbs state at inverse temperature $\beta'$ put $q = e^{-\beta'/r}$; the partition function in the symmetric ordering is
\[
\Phi_{\delta}(q) = q^{B_{2}(\delta)/2}\prod_{w \in W}(1 + q^{w}), \qquad E(\beta') = -\partial_{\beta'}\ln\Phi_{\delta} = \frac{1}{r}\Bigl[\tfrac12 B_{2}(\delta) + \sum_{w \in W}\frac{w}{q^{-w} + 1}\Bigr],
\]
and the entropy is $\ln\Phi_{\delta} + \beta' E$. The thermal test is $E(2\pi) = 0$, where $q = e^{-2\pi/r}$. On the continuum side the boost ground state has null-ordered Killing-frame density $-\kappa^{2}/(48\pi)$ per complex component and the null ground state $0$ (clause (i) of the same theorem), so over a window of Killing time $2\pi r/\kappa$ they carry $-\kappa r/24$ and $0$.

(i) $B_{2}(\delta) = (\delta - \tfrac12)^{2} - \tfrac{1}{12}$, so $\tfrac{1}{2r}B_{2}(\delta) + \tfrac{r}{24} = \tfrac{1}{2r}\bigl[(\delta - \tfrac12)^{2} - \tfrac{1}{12}(1 - r^{2})\bigr]$. The left side of the conic is nonnegative, so $r \le 1$, and $r = 1$ forces $\delta = \tfrac12$. The relabelling $k \mapsto -1 - k$ exchanges the two towers of $W$ and $\delta$ with $1 - \delta$, leaving $W$ unchanged.

(ii) \emph{A lemma.} If $f$ is differentiable on $(0, \infty)$ and $f(c/s) = f(s)$ for a constant $c > 0$, then $f'(s) = -(c/s^{2})f'(c/s)$, so $f'(\sqrt{c}) = 0$. Numbers below are in $\mathrm{span}\{1, I_{4}\}$ and $q = e^{2\pi I_{4}\tau}$ with $\tau = I_{4}\beta'/(2\pi r)$ on the positive imaginary axis, where the functions of Fact \claimref{synthesis.imported-theta-function-identities} are positive.

\emph{Twist $\tfrac12$.} Theorem \claimref{synthesis.antiperiodic-array-against-the-light-ray-continuum-at-window-length-r}(ii): at twist $\tfrac12$ the sea meets the boost ground state's window value, and the symmetric-ordered thermal energy vanishes, exactly when $r = 1$.

\emph{Twist $\tfrac14$.} Theorem \claimref{synthesis.light-ray-continuum-is-the-twist-integral-of-arrays-and-window-values-are-twist-averages}(iii): the stored array is the product of the arrays of twists $\tfrac14$ and $\tfrac34$ at $r = \tfrac12$, each with energy $\pi/24$ above its sea and sea energy $-\pi/24$ in units of $1/\beta$, so its symmetric-ordered thermal energy is $0$, and each carries $\tfrac14\ln 2$. Its sea energy $-\kappa/48$ is the window value $-\kappa r/24$ at $r = \tfrac12$.

\emph{Twist $\tfrac13$.} $B_{2}(\tfrac13) = -\tfrac{1}{18}$ and $W = \{n/3 : n \ge 1, 3 \nmid n\}$. Put $p = q^{1/3} = e^{2\pi I_{4}\sigma}$, $\sigma = \tau/3$. From $\prod_{n \ge 1}(1 + p^{n}) = \prod(1 - p^{2n})/\prod(1 - p^{n}) = p^{-1/24}\eta(2\sigma)/\eta(\sigma)$ and the same identity for $p^{3}$,
\[
\prod_{3 \nmid n}(1 + p^{n}) = \frac{\prod(1 + p^{n})}{\prod(1 + p^{3n})} = p^{1/12}\,\frac{\eta(2\sigma)\eta(3\sigma)}{\eta(\sigma)\eta(6\sigma)}, \qquad \Phi_{1/3} = p^{-1/12}\prod_{3 \nmid n}(1 + p^{n}) = \frac{\eta(2\sigma)\eta(3\sigma)}{\eta(\sigma)\eta(6\sigma)}.
\]
Under $\sigma \mapsto -1/(6\sigma)$, clause (b) of the Fact gives $\eta(2\sigma) \mapsto \eta(-1/(3\sigma)) = \sqrt{-3I_{4}\sigma}\,\eta(3\sigma)$, $\eta(3\sigma) \mapsto \sqrt{-2I_{4}\sigma}\,\eta(2\sigma)$, $\eta(\sigma) \mapsto \sqrt{-6I_{4}\sigma}\,\eta(6\sigma)$ and $\eta(6\sigma) \mapsto \sqrt{-I_{4}\sigma}\,\eta(\sigma)$; the square roots are positive on the axis and $\sqrt{3}\sqrt{2} = \sqrt{6}\sqrt{1}$, so $\Phi_{1/3}$ is invariant. With $\sigma = I_{4}s$ the lemma ($c = \tfrac16$) gives the stationary point $s = 1/\sqrt{6}$; since $s = \beta'/(6\pi r)$, at $\beta' = 2\pi$ this is $r = \sqrt{6}/3 = \sqrt{2/3}$.

In each case the point lies on the conic of (i): $(1 - r^{2})/12 = 0, \tfrac{1}{16}, \tfrac{1}{36}$ is $(\delta - \tfrac12)^{2}$. At a point passing both tests the excitation energy is minus the sea energy, $\kappa r/24$. The mirror twists follow from (i). The check encloses the thermal test value at the three points in rational intervals of width below $10^{-12}$ containing $0$, and excludes $0$ at five control points.

(iii) Where $E(2\pi) = 0$ the entropy is $\ln\Phi_{\delta}$. At twist $\tfrac12$ it is $\tfrac12\ln 2$ (Theorem \claimref{synthesis.half-integer-weight-array-carries-half-ln2}(i)); at twist $\tfrac14$, $r = \tfrac12$, it is $\tfrac14\ln 2$ by the cited product. At twist $\tfrac13$ the rigorous rational bounds of the check give $\sum_{w}s(2\pi w/r) \in (0.29379, 0.29380)$ with $s(x) = \ln(1 + e^{-x}) + x/(e^{x} + 1)$; no closed form is claimed. Dividing by $r$ gives the values per thermal period; the bounds $3S^{2}/2 \in (0.3598^{2}, 0.3599^{2})$, $\tfrac12\ln 2 < 0.3466$ and $\pi/6 > 0.5235$ are in the check. The continuum's rate $\pi/6$ per thermal period per complex component is the extensive value $\tfrac{1}{2\pi}\int_{\R}s(x)\,dx = \pi/6$ (the integral is $\pi^{2}/3$), a standard evaluation quoted for comparison and used in no step.

(iv) The Killing flow at one modular period multiplies the mode $k$ by the phase of angle $\beta\omega_{k} = 2\pi(k + \delta)/r$ in the plane of the boundary modes' unit $J$; it is $\psi \mapsto -\psi$ exactly when every such angle is an odd multiple of $\pi$. Differences of consecutive angles give $1/r = m \in \Z$, and then $2m\delta$ is an odd integer $2j + 1$. On these, $(\delta - \tfrac12)^{2} = p^{2}/(4m^{2})$ with $p = 2j + 1 - m$, and the conic reads $3p^{2} = m^{2} - 1$. For $m = 1$: $p = 0$, the stored array. For $m = 2$: $p = \pm 1$, twists $\tfrac14$, $\tfrac34$, which pass by (ii). For $3 \le m \le 6$ there is no solution. For $m \ge 7$ on the conic, every $w \in W$ exceeds $\tfrac15$ because $|\delta - \tfrac12| \le 1/\sqrt{12} < \tfrac{3}{10}$; since $ue^{-2\pi mu}$ decreases for $u > 1/(2\pi m)$,
\[
\sum_{w \in W}\frac{w}{e^{2\pi mw} + 1} < 2\sum_{k \ge 0}(k + \tfrac15)e^{-2\pi m(k + 1/5)} < 3e^{-2\pi m/5} < \frac{1}{24m^{2}} = -\tfrac12 B_{2}(\delta),
\]
the last inequality holding at $m = 7$ and then for all larger $m$ because $m^{2}e^{-2\pi m/5}$ decreases ($(1 + 1/m)^{2} \le \tfrac94 < e^{2\pi/5}$). So the thermal test value is negative. The twist $\tfrac13$ has $1/r = \sqrt{3/2}$, not an integer. Particle-hole symmetry of one array is $\delta = 1 - \delta$, and on the conic $\delta = \tfrac12$ forces $r = 1$; the spacing $\kappa$ is $r = 1$ directly.

\emph{Closing sentence and remarks.} With affine weights $ck + d$ the slope of Theorem \claimref{synthesis.q-splits-maximum-entropy-affine-weights}(iv) is $y = \beta c = 2\pi/r$, so the three points are $y = 2\pi$, $4\pi$, $2\pi\sqrt{3/2}$; that theorem records the offset class and $y = 2\pi$ as the two choices of (Q). The entropies per thermal period of (iii) differ only at the third point. For Remark 2: at twist $\tfrac13$ the modes $k = 0, 1, -2$ have weights proportional to $\tfrac13$, $\tfrac43$, $-\tfrac53$, whose sum is $0$, a signed sum of three weights, so clause 3 of Theorem \claimref{synthesis.univalence-from-even-stationary-dynamics} (applied to any finite set of modes containing these three) gives stationary states that are not even; at twists $\tfrac12$ and $\tfrac14$ the weights are odd multiples of $\tfrac{1}{2r}$ and $\tfrac{1}{4r}$, and a signed sum of an odd number of odd integers is odd, hence nonzero. The Gibbs state is even by clause 2 there.

\emph{What is not shown.} That no further point of the conic passes the thermal test when the parity clause is dropped: a floating-point scan (step $0.005$ in $r$, with bisection at the sign changes) finds the three points only, which is a computation and not a proof. That the tests assume no reading is argued at clause level, not audited over the cone, which contains (Q), (N$_{+}$), J3 and the window reading. Nothing here derives the spacing, and nothing is claimed about which array the count should use.
